\section{消没定理}\srcpage{13}

设 $i\ge 1$ 与 $d\ge 0$ 为整数。考虑图表
\[
\begin{tikzcd}
 & \Mod^d_i(X) \arrow[dl,"p'"] \arrow[r,"q"] \arrow[dr,"p"] & \Coh^d_0(X) \\
\Fib_i(X) & & \Fib_i(X)
\end{tikzcd}
\]

其中 $\Mod^d_i(X)$ 是由对 $(\cL,\cL'\subset\cL)$ 构成的叠。这里 $\cL$ 是秩 $i$ 的向量丛，$\cL'\subset\cL$ 是余长 $d$ 的下修改。投影 $p'$ 与 $p$ 分别把 $(\cL,\cL'\subset\cL)$ 送到 $\cL'$ 与 $\cL$。$q$ 则把它送到凝聚 $\cO_X$-模 $\cL/\cL'$，后者一般秩为 $0$、长度为 $d$。

对 $X$ 上任一局部系统 $E$，第 4 节的函子 $M'^d_{i,E}$ 有如下变体
\[
M^d_{i,E}:\Dbc(\Fib_i(X))\to\Dbc(\Fib_i(X)),\qquad
M^d_{i,E}(K)=p_*p'^*(K\otimes q^*\cL^d_E)[id],
\]
\srcpage{14}
其中 $\cL^d_E$ 是 $\Coh^d_0(X)$ 上由 3.2 定义的反常层。

\begin{theorem}[Gaitsgory, {[Ga]}]
设 $E$ 是秩 $n$ 的不可约局部系统。则对一切 $i=1,\dots,n-1$ 与一切整数 $d>in(2g-2)$\yizhu{原文作 $in(2g-2)$，其中 $in$ 是 $i$ 与 $n$ 之积，即 $d>i\,n(2g-2)$；注 6.2 中 $i=1$ 的情形正写作 $n(2g-2)$，可互证。}{术语背景，A（原文注 6.2 互证）}，函子 $M^d_{i,E}$ 恒为零。
\end{theorem}

\begin{remark}
这个定理的证明很长，技术性很强。它远远超出了本报告的范围。然而我们指出，当 $i=1$ 时，这一陈述不过是 Deligne 一个结果的重新表述（见 [Dr 2]），该结果说
\[
(\varphi^{(d)})_*E^{(d)}=0
\]
对 $X$ 上任一秩 $n>1$ 的不可约局部系统与任一整数 $d>n(2g-2)$ 成立，其中 $\varphi^{(d)}:X^{(d)}\to\Pic^0(X)$ 是第 1 节中考虑的射影丛。
\end{remark}

Frenkel、Gaitsgory 与 Vilonen 用定理 6.1 证明下述命题。

\begin{proposition}
设 $E$ 是秩 $n$ 的不可约局部系统。则对每个 $i=1,\dots,n-1$，典范态射
\[
(j_i)_!K_{E,i}\to(j_i)_*K_{E,i}
\]
是同构。
\end{proposition}

证明思路：在 $\cC_i$ 的开集 $\Fib_i(X)\cap\cC_i$ 上，$\cV^\circ_i\subset\cV_i$ 是向量丛 $\cV_i\to\cC_i$ 中零截面的补集。由引理 (0.3)，要证明该开集上的断言，只需证明 $(\overline{\pi}^\circ_i)_!K_{E,i}=(0)$，其中 $\overline{\pi}^\circ_i:\overline{\cV}^\circ_i\to\cC_i$ 是投影；或等价地，证明关于投影 $\Pv(\cV_i)\to\cC_i$ 的同一等式，因为这两个投影在 $\Fib_i(X)\cap\cC_i$ 上重合。而由引理 5.2，对每个整数 $d\ge 0$，有
\[
(\overline{\pi}^\circ_i)_!K_{E,i}|(\cC_i\cap\Fib_i^{i(i-1)(g-1)+d}(X))
=M^d_{i,E}(\pi'_*L_i)|(\cC_i\cap\Fib_i^{i(i-1)(g-1)+d}(X))
\]
其中 $\pi':\Fib'_i(X)\to\Fib_i(X)$ 是遗忘截面 $s:(\OmX)^{\otimes n-1}\to\cL$ 的态射，于是由定理 6.1 得结论。

由命题 6.3，若 $K_{E,i}$ 是不可约反常层，则 $(j_i)_!K_{E,i}$ 也是，因为此时它是 $K_{E,i}$ 的中间延拓。由于齐次 Fourier 变换保持反常性与不可约性，对 $i$ 作归纳可知：$K_{E,i}$ 在 $\overline{\cV}^\circ_i$ 的每个连通分支上的限制是不可约反常层。特别地，我们已经证明：

\srcpage{15}
\begin{corollary}
设 $E$ 是秩 $n$ 的不可约局部系统。则 $K_{E,n}$ 在 $\overline{\cV}^\circ_n$ 的每个连通分支上的限制是不可约反常层。
\end{corollary}
